\scanpage{48}\mspage{41}
% archive sheet number: 47
On peut songer à utiliser une factorisation \uline{fonctorielle} \unsure{des} flèches de $M$ en
\[
\begin{aligned}
f=pi\quad &\text{avec}\quad i\in TC,\ p\in F\\
&\text{ou}\quad i\in C,\ p\in TF,
\end{aligned}
\]
pour déduire une factorisation \unsure{simultanée} pour les flèches $\unsure{f.|f_{\cdot}}$ de $M(I)$ en
\[
\begin{aligned}
f=pi\qquad &i\in TC(I),\ p\in F(I)\\
&\text{ou}\quad i\in C(I),\ p\in TF(I),
\end{aligned}
\]
où $C(I),TC(I),F(I),TF(I)$ sont ici définis via $C,TC,F,TF$ ``argument par argument''. Mais l'ennui, c'est que en procédant de cette façon naïve, on n'aura plus nécessairement $C(I)\longleftrightarrow TF(I)$, $TC(I)\unsure{\longleftrightarrow|\longleftarrow} F(I)$. Il faut considérablement rapetisser $TF(I)$ et $F(I)$, pour avoir une propriété RLP \unsure{p.r.à|p.r.} $C(I)$ resp. $TC(I)$.

Donc je ne vois guère d'autre moyen de s'en tirer, dans ce contexte général, qu'en supposant la stabilité de $W$ par $\unsure{\varinjlim|\underline{\lim}}$ filtrantes (ce qui n'est après \unsure{tout} pas très méchant).
